\documentclass[11pt,a4paper]{article}

\usepackage[margin=2.5cm]{geometry}
\usepackage{amsmath,amssymb,amsthm}
\usepackage{booktabs}
\usepackage{graphicx}
\usepackage[most]{tcolorbox}
\usepackage{float}

\newtheorem{theorem}{Theorem}[section]
\newtheorem{proposition}[theorem]{Proposition}
\newtheorem{conjecture}[theorem]{Conjecture}
\theoremstyle{definition}
\newtheorem{definition}[theorem]{Definition}

\newcommand{\R}{\mathbb{R}}
\newcommand{\Z}{\mathbb{Z}}
\newcommand{\Q}{\mathbb{Q}}
\newcommand{\N}{\mathbb{N}}
\newcommand{\G}{\Gamma}

\title{\textbf{Approximating Orthogonal Functions\\ with Orthogonal Polynomials}}
\author{Robin Langer and Claude Code}
\date{August 2026}

\begin{document}
\maketitle

\begin{abstract}
We can always approximate a function by a polynomial, and the
approximation gets better as the degree of the polynomial goes up.
The Hermite functions --- the eigenfunctions of the Fourier
transform --- can be approximated by Hermite polynomials. The
approximation coefficients form a triangular matrix with integer
entries $D_{mk}$, which we conjecture to count weighted bicolour
matchings. For the Chebyshev polynomials only two matchings survive,
and the matrix is built from quartic factorials. We extend this
construction to the Meixner--Pollaczek
polynomials at $\lambda = \tfrac14$ and $\tfrac34$. In this case
there are constants $A$ and $B$ such that every coefficient in the
approximation matrix is of the form $sA + tB$ where $s$ and $t$ are
integers. On the critical line, the Mellin transform carries the
Hermite functions unitarily onto these two families.
\end{abstract}

%% ===================================================================
\section{The Construction}
%% ===================================================================

\begin{tcolorbox}[breakable, colback=black!5, colframe=black!5,
  boxrule=0pt, sharp corners]
Each family follows the same pattern, in three steps.

\textbf{Step 1: the polynomials.} Fix a kernel $W \ge 0$ on a set $X$.
It defines an inner product on polynomials,
\[
  \langle f, g \rangle_W = \int_X f\,g\;W.
\]

Gram--Schmidt produces the orthogonal polynomials $P_n$:
\[
  \langle P_m, P_n \rangle_W = h_n \delta_{mn}.
\]

Normalise:
\[
  \widehat P_n = P_n/\sqrt{h_n}.
\]

\medskip
\textbf{Step 2: the functions.} Multiplying by $\sqrt W$ turns them
into the \emph{associated functions}
\[
  \Phi_n = \widehat P_n\, \sqrt W .
\]

These are orthonormal for the plain inner product:
\[
  \int_X \Phi_m\, \Phi_n = \delta_{mn}.
\]

When $\sqrt W$ is not analytic, we instead split $W = w_+ w_-$
analytically and take the \emph{dual pair of functions}
\[
  F_n = w_+\, \widehat P_n, \qquad G_n = w_-\, \widehat P_n,
  \qquad \int_X F_m\, G_n = \delta_{mn}.
\]

\textbf{Step 3: the coefficients.} We seek \emph{polynomial
approximations of the functions}. The best degree-$N$ approximation
of $\Phi_m$ in $L^2(W)$ is given by
\[
  \sum_{k \le N} b_{mk}\, \widehat P_k .
\]

The coefficients give a \textbf{triangular matrix}.

\medskip
The coefficients are given by
\[
  b_{mk} = \langle \Phi_m, \widehat P_k \rangle_W
         = \int_X \bigl(\widehat P_m \sqrt W\bigr)\, \widehat P_k\; W
         = \int_X \widehat P_m\, \widehat P_k\; W^{3/2}.
\]

In the dual case there are two coefficient matrices, one for each
function of the pair:
\[
  c^+_{mk} = \langle F_m, \widehat P_k \rangle_W
         = \int_X \widehat P_m\, \widehat P_k\; w_+ W
  \qquad (w_+ W = w_+^2\, w_-),
\]
\[
  c^-_{mk} = \langle G_m, \widehat P_k \rangle_W
         = \int_X \widehat P_m\, \widehat P_k\; w_- W
  \qquad (w_- W = w_-^2\, w_+).
\]
\end{tcolorbox}

%% ===================================================================
\section{Families of Orthogonal Polynomials}
%% ===================================================================

%% -------------------------------------------------------------------
\subsection{Hermite}

In the Hermite polynomial case the kernel is given by
$W = e^{-x^2}$ on $\R$, with $h_n = 2^n n! \sqrt\pi$.
The Hermite functions are given by $\Phi_n = \widehat H_n e^{-x^2/2}$.
We are interested in the coefficients of the polynomial approximation
of the Hermite function by the Hermite polynomials
\[
  b_{mk} = \int \widehat H_m \widehat H_k\, e^{-3x^2/2} dx .
\]
With $b_{00} = \sqrt{2/3}$, define the rescaled integers
\[
  D_{mk} \;=\; \frac{b_{mk}}{b_{00}}\;
  \sqrt{2^m m!\; 2^k k!}\;\; 3^{(m+k)/2} \;\in\; \Z .
\]
$D_{mk} = 0$ for $m+k$ odd. The even entries ($m \le k \le 7$):

\begin{tcolorbox}[colback=black!5, colframe=black!5, boxrule=0pt, sharp corners]
\begin{center}
\begin{tabular}{r|rrrr}
$D_{mk}$ & $k=m$ & $m{+}2$ & $m{+}4$ & $m{+}6$ \\
\hline
$m=0$ & $1$      & $-2$      & $12$      & $-120$ \\
$1$   & $4$      & $-24$     & $240$     & $-3360$ \\
$2$   & $36$     & $-408$    & $6000$    &  \\
$3$   & $528$    & $-9120$   & $181440$  &  \\
$4$   & $10896$  & $-254880$ &           &  \\
$5$   & $290880$ & $-8588160$&           &  \\
$6$   & $9529920$&           &           &  \\
$7$   & $369895680$ &        &           &  \\
\end{tabular}
\end{center}
\end{tcolorbox}

\begin{tcolorbox}[colback=black!5, colframe=black!5, boxrule=0pt, sharp corners]
\begin{theorem}[Hermite integrality]\label{thm:hermite-int}
$D_{mk} \in \Z$, vanishing for $m+k$ odd. Explicitly, with
$b = \tfrac{m-a}{2}$, $c = \tfrac{k-a}{2}$,
\[
  D_{mk} \;=\; \sum_{a}\; 4^a\, (-2)^{b+c}\, a!\;
  \frac{m!}{a!\, b!\, 2^b}\; \frac{k!}{a!\, c!\, 2^c}\,,
\]
and each factor $\dfrac{n!}{a!\, b!\, 2^b}$ is an integer: it counts
the partial matchings of $[n]$ with $a$ unmatched points and $b$ pairs.
\end{theorem}
\end{tcolorbox}

\begin{tcolorbox}[colback=black!5, colframe=black!5, boxrule=0pt, sharp corners]
\begin{conjecture}[Hermite]\label{thm:hermite}
$D_{mk}$ is a weighted count of perfect matchings. Colour $m$ points
red and $k$ points blue; then
\[
  D_{mk} \;=\; \sum_{\substack{\text{perfect matchings}\\
  \text{of the } m+k \text{ points}}}
  4^{\#\{\text{red--blue pairs}\}} \; (-2)^{\#\{\text{same-colour pairs}\}} .
\]
\end{conjecture}
\end{tcolorbox}

\begin{tcolorbox}[colback=black!5, colframe=black!5, boxrule=0pt, sharp corners]
\begin{theorem}[Hermite diagonal]\label{thm:hermite-diag}
The diagonal is manifestly positive:
\[
  \frac{D_{nn}}{n!}
  \;=\; \sum_{j} \binom{n}{2j} \binom{2j}{j}\, 4^{\,n-2j}
  \;=\; 1, 4, 18, 88, 454, 2424, \ldots \quad (\text{OEIS A081671}),
\]
the largest coefficient of $(1 + 4x + x^2)^n$, with ordinary
generating function $1/\sqrt{(1-2x)(1-6x)}$ and exponential
generating function $e^{4x} I_0(2x)$.
\end{theorem}
\end{tcolorbox}


\subsubsection*{Reverse}

The \emph{reverse} expansion $\widehat H_m = \sum_k \widetilde b_{mk}\, \Phi_k$
has one power of $\sqrt W$ in place of three:
\[
  \widetilde b_{mk} = \int \widehat H_m \widehat H_k\, e^{-x^2/2}\, dx .
\]
With $\widetilde b_{00} = \sqrt2$, the same rescaling without the
power of $3$,
\[
  S_{mk} \;=\; \frac{\widetilde b_{mk}}{\widetilde b_{00}}\,
  \sqrt{2^m m!\; 2^k k!} \;\in\; \N .
\]
$S_{mk} = 0$ for $m+k$ odd. The even entries ($m \le k \le 7$):

\begin{tcolorbox}[colback=black!5, colframe=black!5, boxrule=0pt, sharp corners]
\begin{center}
\begin{tabular}{r|rrrr}
$S_{mk}$ & $k=m$ & $m{+}2$ & $m{+}4$ & $m{+}6$ \\
\hline
$m=0$ & $1$      & $2$      & $12$      & $120$ \\
$1$   & $4$      & $24$     & $240$     & $3360$ \\
$2$   & $36$     & $408$    & $6000$    &  \\
$3$   & $528$    & $9120$   & $181440$  &  \\
$4$   & $10896$  & $254880$ &           &  \\
$5$   & $290880$ & $8588160$&           &  \\
$6$   & $9529920$&           &           &  \\
$7$   & $369895680$ &        &           &  \\
\end{tabular}
\end{center}
\end{tcolorbox}

\begin{tcolorbox}[colback=black!5, colframe=black!5, boxrule=0pt, sharp corners]
\begin{theorem}[Hermite, reverse direction]\label{thm:hermite-rev}
$S_{mk} = (-1)^{(k-m)/2}\, D_{mk}$ --- the forward matrix with the
signs stripped. In the matchings of Conjecture~\ref{thm:hermite} the
same-colour weight $-2$ is replaced by $+2$, making the sum
manifestly positive:
\[
  S_{mk} \;=\; \sum_{\substack{\text{perfect matchings}\\
  \text{of the } m+k \text{ points}}}
  4^{\#\{\text{red--blue pairs}\}} \; 2^{\#\{\text{same-colour pairs}\}} .
\]
\end{theorem}
\end{tcolorbox}

%% -------------------------------------------------------------------
\subsection{Chebyshev}

In the Chebyshev case the kernel is $W = (1-x^2)^{-1/2}$ on $(-1,1)$,
with $h_0 = \pi$ and $h_n = \pi/2$. The associated functions are
$\Phi_n = \widehat T_n\, (1-x^2)^{-1/4}$. Substituting $x = \cos\theta$
turns every matching integral into a two-term sum,
\[
  \int_{-1}^{1} T_m T_k\, (1-x^2)^{-3/4}\, dx
  \;=\; \tfrac12\bigl[\, J(m-k) + J(m+k) \,\bigr],
  \qquad
  J(q) = \int_0^\pi \frac{\cos q\theta}{\sqrt{\sin\theta}}\, d\theta ,
\]
because $T_m T_k$ linearises into two cosines. Where Hermite sums over
all matchings, Chebyshev keeps only two.

Let $P(n) = 1 \cdot 5 \cdots (4n-3)$ and $Q(n) = 3 \cdot 7 \cdots
(4n-1)$ denote the quartic factorials. Then
\[
  J(0) = \frac{\sqrt\pi\; \G(\tfrac14)}{\G(\tfrac34)},
  \qquad
  J(2n) = J(0)\, \frac{P(n)}{Q(n)} \;>\; 0 .
\]
Every coefficient is a rational multiple of the single constant
$J(0)$. With $s = \tfrac{m+k}{2}$, $d = \tfrac{|m-k|}{2}$, define
\[
  D_{mk} \;=\; \frac{Q(s)}{J(0)}
  \int_{-1}^{1} T_m\, T_k\, (1-x^2)^{-3/4}\, dx .
\]

\begin{tcolorbox}[colback=black!5, colframe=black!5, boxrule=0pt, sharp corners]
\begin{theorem}[Chebyshev integrality]\label{thm:cheb-int}
$D_{mk} = 0$ for $m+k$ odd, and otherwise
\[
  D_{mk} \;=\; \frac12 \left[\, P(d)\,\frac{Q(s)}{Q(d)} + P(s) \,\right]
\]
is a \emph{positive} integer: both bracketed terms are products of odd
numbers, so their sum is even.
\end{theorem}
\end{tcolorbox}

The entries ($m \le k \le 7$):

\begin{tcolorbox}[colback=black!5, colframe=black!5, boxrule=0pt, sharp corners]
\begin{center}
\begin{tabular}{r|rrrr}
$D_{mk}$ & $k=m$ & $m{+}2$ & $m{+}4$ & $m{+}6$ \\
\hline
$m=0$ & $1$        & $1$         & $5$          & $45$ \\
$1$   & $2$        & $6$         & $50$         & $630$ \\
$2$   & $13$       & $61$        & $705$        & $11385$ \\
$3$   & $138$      & $870$       & $12810$      & $251910$ \\
$4$   & $2025$     & $15945$     & $284685$     & $6592725$ \\
$5$   & $37890$    & $356790$    & $7477650$    & $199153350$ \\
$6$   & $861525$   & $9424485$   & $226586025$  & $6818954625$ \\
$7$   & $23052330$ & $286937910$ & $7779098250$ & $260942613750$ \\
\end{tabular}
\end{center}
\end{tcolorbox}

\begin{tcolorbox}[colback=black!5, colframe=black!5, boxrule=0pt, sharp corners]
\begin{theorem}[Chebyshev diagonal and edge]\label{thm:cheb-diag}
The edge is the quartic factorial, $D_{0,2n} = P(n)$
(OEIS A007696), and the diagonal is the \emph{average of the two
quartic factorials},
\[
  D_{nn} = \frac{Q(n) + P(n)}{2}
  \;=\; 1,\, 2,\, 13,\, 138,\, 2025,\, 37890,\, 861525,\, \ldots
\]
This sequence is not in the OEIS.
\end{theorem}
\end{tcolorbox}

\subsubsection*{Reverse}

The reverse expansion uses $W^{1/2} = (1-x^2)^{-1/4}$ in place of
$W^{3/2}$, hence $\widetilde J(q) = \int_0^\pi \cos q\theta\,
\sqrt{\sin\theta}\; d\theta$. The third quartic class
$\widetilde P(n) = 5 \cdot 9 \cdots (4n+1)$ appears, and the ratios
turn negative,
\[
  \widetilde J(0) = \frac{\sqrt\pi\;\G(\tfrac34)}{\G(\tfrac54)},
  \qquad
  \widetilde J(2n) = -\,\widetilde J(0)\,
  \frac{Q(n-1)}{\widetilde P(n)} \quad (n \ge 1).
\]
Define $S_{mk} = \dfrac{\widetilde P(s)}{\widetilde J(0)}
\displaystyle\int_{-1}^{1} T_m T_k\, (1-x^2)^{-1/4} dx$. Then
$S_{mk} \in \Z$ by the same parity argument, with edge
$S_{0,2n} = -Q(n-1)$ and diagonal
$S_{nn} = \tfrac12 [\widetilde P(n) - Q(n-1)] = 1, 2, 21, 282, 4857,
\ldots$ (also not in the OEIS). The entries ($m \le k \le 7$):

\begin{tcolorbox}[colback=black!5, colframe=black!5, boxrule=0pt, sharp corners]
\begin{center}
\begin{tabular}{r|rrrr}
$S_{mk}$ & $k=m$ & $m{+}2$ & $m{+}4$ & $m{+}6$ \\
\hline
$m=0$ & $1$        & $-1$         & $-3$          & $-21$ \\
$1$   & $2$        & $-6$         & $-30$         & $-294$ \\
$2$   & $21$       & $-69$        & $-447$        & $-5481$ \\
$3$   & $282$      & $-1110$      & $-8694$       & $-126630$ \\
$4$   & $4857$     & $-22617$     & $-206955$     & $-3474765$ \\
$5$   & $102690$   & $-555030$    & $-5804190$    & $-110124630$ \\
$6$   & $2577645$  & $-15898365$  & $-186995655$  & $-3951960705$ \\
$7$   & $74949210$ & $-520103430$ & $-6796188630$ & $-158228220150$ \\
\end{tabular}
\end{center}
\end{tcolorbox}

The sign pattern is inverted relative to Hermite: there the
approximation matrix $D$ carried the signs and the reverse matrix $S$
was positive; here $D$ is positive and $S$ carries the signs. The
three quartic residue classes share the work --- $P$ and $Q$ build the
forward matrix, $Q$ and $\widetilde P$ the reverse.

\subsubsection*{The U-kernel}

The kernel $W = (1-x^2)^{1/2}$ (Chebyshev-U) uses the fourth and
last quartic class $Q^{+}$. Here $U_m U_k$ linearises into a
\emph{difference} of cosines, and with
$K(q) = \int_0^\pi \cos q\theta\, \sin^{3/2}\theta\, d\theta$,
\[
  Q^{+}(n) = 7 \cdot 11 \cdots (4n+3),
  \qquad
  K(0) = \frac{J(0)}{3},
  \qquad
  K(2n) = -\,3K(0)\, \frac{P(n-1)}{Q^{+}(n)} \quad (n \ge 1).
\]
Define
\[
  D^U_{mk} \;=\; \frac{Q^{+}(s+1)}{K(0)}
  \int_0^\pi \sin\bigl((m{+}1)\theta\bigr) \sin\bigl((k{+}1)\theta\bigr)
  \sqrt{\sin\theta}\; d\theta .
\]

\begin{tcolorbox}[colback=black!5, colframe=black!5, boxrule=0pt, sharp corners]
\begin{theorem}[U integrality]\label{thm:chebU-int}
$D^U_{mk} = 0$ for $m+k$ odd, and otherwise
\[
  D^U_{mk} \;=\; \frac12 \left[\, Q^{+}(s{+}1)\,\rho(d) + 3P(s) \,\right]
  \;\in\; \Z,
  \qquad
  \rho(0)=1, \quad \rho(n) = -\,\frac{3P(n-1)}{Q^{+}(n)} ,
\]
with edge $D^U_{0,2n} = -15\,P(n-1)$ for $n \ge 1$ (the quartic
factorial A007696 again) and diagonal
\[
  D^U_{nn} = \frac{Q^{+}(n{+}1) + 3P(n)}{2}
  \;=\; 5,\, 40,\, 585,\, 11040,\, 253245,\, \ldots
\]
not in the OEIS. The diagonal is positive and every other entry
negative --- a third sign pattern, after the all-positive $D$ and the
alternating $S$.
\end{theorem}
\end{tcolorbox}

\subsubsection*{An asymmetric kernel: Catalan numbers}

The kernel $W = (1-x)^{1/2}(1+x)^{-1/2}$ breaks the $x \mapsto -x$
symmetry; its orthogonal polynomials are the fourth-kind Chebyshev
$W_n(\cos\theta) = \sin\bigl((n{+}\tfrac12)\theta\bigr)/\sin(\theta/2)$,
with $h_n = \pi$ uniformly. The Wallis power drops to zero, the
constant $\G(\tfrac14)$ disappears, and the whole matching matrix lies
in $\Z[\tfrac12]\cdot\sqrt2$ --- dyadic rationals. The first row is
classical combinatorics: with $b_{0k} = \frac1\pi \int W_0 W_k
W^{3/2}\, dx$,
\[
  2^k \sqrt2\; b_{0k} \;=\;
  \begin{cases}
    2\binom{2n}{n} & k = 2n \ge 2, \\[4pt]
    -\,(4n+3)\,\mathrm{Cat}(n) & k = 2n+1,
  \end{cases}
  \qquad
  2^0\sqrt2\; b_{00} = 3,
\]
giving $3,\, -3,\, 4,\, -7,\, 12,\, -22,\, 40,\, -75,\, 140,\, -266,
\ldots$ (verified exactly through $k=13$; the signed sequence is not
in the OEIS). The quartic weights of the symmetric kernels cannot
count lattice paths; the asymmetric kernel apparently can.

In total the Chebyshev-type matrices use the four quartic
factorial classes $1{\cdot}5{\cdots}$, $3{\cdot}7{\cdots}$,
$5{\cdot}9{\cdots}$, $7{\cdot}11{\cdots}$ --- the four odd
arithmetic progressions of difference four, each exactly once.

%% -------------------------------------------------------------------
\subsection{Laguerre}

For the Laguerre polynomials with $\alpha = 0$ the kernel is
$W = e^{-x}$ on $(0,\infty)$. We denote the Laguerre polynomials by
$L_n$ and the Laguerre functions by
$\Phi_n = L_n e^{-x/2}$. They are already orthonormal ($h_n = 1$). The coefficients
\[
  b_{mk} = \int_0^\infty L_m L_k\, e^{-3x/2} dx ,
\]
are all exactly rational. We rescale:
\[
  b_{mk} \;=\; \frac{2\, E_{mk}}{3^{\,m+k+1}}, \qquad E_{mk} \in \N :
\]

\begin{tcolorbox}[colback=black!5, colframe=black!5, boxrule=0pt, sharp corners]
\begin{center}
\begin{tabular}{r|rrrrrr}
$E_{mk}$ & $k=0$ & $1$ & $2$ & $3$ & $4$ & $5$ \\
\hline
$m=0$ & $1$ & $1$ & $1$ & $1$ & $1$ & $1$ \\
$1$ &  & $5$ & $9$ & $13$ & $17$ & $21$ \\
$2$ &  &  & $33$ & $73$ & $129$ & $201$ \\
$3$ &  &  &  & $245$ & $593$ & $1181$ \\
$4$ &  &  &  &  & $1921$ & $4881$ \\
$5$ &  &  &  &  &  & $15525$ \\
\end{tabular}
\end{center}
\end{tcolorbox}

\begin{tcolorbox}[colback=black!5, colframe=black!5, boxrule=0pt, sharp corners]
\begin{theorem}[Laguerre]\label{thm:laguerre}
$E_{mk}$ counts weighted Delannoy paths:
\[
  E_{mk} \;=\; \sum_{\text{paths } (0,0) \to (m,k)} 3^{\#\text{diagonal steps}}
  \;=\; \sum_{j \ge 0} \binom{m}{j} \binom{m+k-j}{m}\, 3^{\,j},
\]
the sum over monotone lattice paths with steps $(1,0)$, $(0,1)$,
$(1,1)$, a diagonal step weighing $3$. In particular
$E_{1,k} = 4k+1$ and
$E_{nn} = \sum_j \binom{n}{j}^2 4^j = {}_2F_1(-n,-n;1;4)$.
\end{theorem}
\end{tcolorbox}

\subsubsection*{Reverse}

The \emph{reverse} expansion $L_m = \sum_k \widetilde b_{mk}\, \Phi_k$
gives the same Delannoy numbers:

\begin{tcolorbox}[colback=black!5, colframe=black!5, boxrule=0pt, sharp corners]
\begin{center}
\begin{tabular}{r|rrrrrr}
$\widetilde b_{mk}$ & $k=0$ & $1$ & $2$ & $3$ & $4$ & $5$ \\
\hline
$m=0$ & $2$ & $-2$ & $2$ & $-2$ & $2$ & $-2$ \\
$1$ &  & $10$ & $-18$ & $26$ & $-34$ & $42$ \\
$2$ &  &  & $66$ & $-146$ & $258$ & $-402$ \\
$3$ &  &  &  & $490$ & $-1186$ & $2362$ \\
$4$ &  &  &  &  & $3842$ & $-9762$ \\
$5$ &  &  &  &  &  & $31050$ \\
\end{tabular}
\end{center}
\end{tcolorbox}

\begin{tcolorbox}[colback=black!5, colframe=black!5, boxrule=0pt, sharp corners]
\begin{theorem}[Laguerre, reverse direction]\label{thm:laguerre-rev}
\[
  \widetilde b_{mk} = \int_0^\infty L_m L_k\, e^{-x/2}\, dx
  \;=\; 2\,(-1)^{m+k}\, E_{mk} .
\]
Reversing the direction strips the power of $3$ and leaves a sign.
\end{theorem}
\end{tcolorbox}


%% ===================================================================
\section{Not quite combinatorial}
%% ===================================================================

%% -------------------------------------------------------------------
\subsection{Meixner--Pollaczek ($\lambda = \tfrac14$): a dual pair}

For the Meixner--Pollaczek polynomials with $\lambda = \tfrac14$ the
kernel is $W = \G(\tfrac14+ix)\,\G(\tfrac14-ix)$ on $\R$, with
$h_n = h_0 \binom{2n}{n}/4^n$, $h_0 = \sqrt2\, \pi^{3/2}$.
We denote the Meixner--Pollaczek polynomials by $P_n$.
Here $\sqrt W$ has branch points at
$\pm\tfrac{i}{4}$; the analytic split is $w_\pm = \G(\tfrac14 \pm ix)$
and the dual pair of functions is given by
\[
  F_n = \G(\tfrac14+ix)\,\widehat P_n, \qquad
  G_n = \G(\tfrac14-ix)\,\widehat P_n, \qquad
  \int_\R F_m G_n\, dx = \delta_{mn}.
\]
The coefficients of the approximation matrix are given by:
\[
  c^+_{mk} = \int \widehat P_m \widehat P_k\,
  \G(\tfrac14+ix)^2\,\G(\tfrac14-ix)\, dx, \qquad
  c^-_{mk} = \int \widehat P_m \widehat P_k\,
  \G(\tfrac14+ix)\,\G(\tfrac14-ix)^2\, dx
\]
They are complex conjugates of each other, $c^-_{mk} = \overline{c^+_{mk}}$.
With the rescaling
\[
  T_{mk} = \frac{2^m m!\; 2^k k!}{h_0}
  \int_\R P_m P_k\, \G(\tfrac14+ix)^2 \G(\tfrac14-ix)\, dx,
  \qquad
  c^+_{mk} = \frac{T_{mk}}{m!\,k!\,
  \sqrt{\smash[b]{\tbinom{2m}{m}\tbinom{2k}{k}}}} .
\]
Up to a factor of $i$ when $m+k$ is odd, every entry lies in
$\Z A + \Z B$, where
\[
  A = \sqrt{2e}\,K_0(\tfrac12), \qquad
  B = \sqrt{2e}\,K_1(\tfrac12),
\]
and the $K_\nu$ are the modified Bessel functions of the third kind
\[
  K_\nu(z) = \int_0^\infty e^{-z\cosh t} \cosh(\nu t)\, dt :
\]

\begin{tcolorbox}[colback=black!5, colframe=black!5, boxrule=0pt, sharp corners]
\begin{center}
\begin{tabular}{cl@{\qquad}cl}
\toprule
$(m,k)$ & $T_{mk}$ & $(m,k)$ & $T_{mk}$ \\
\midrule
$(0,0)$ & $A$                   & $(1,3)$ & $2575A - 1440B$ \\
$(0,1)$ & $3A - 2B$             & $(1,4)$ & $39411A - 21978B$ \\
$(0,2)$ & $12B - 23A$           & $(1,5)$ & $394932B - 707389A$ \\
$(0,3)$ & $138B - 241A$         & $(2,2)$ & $2369A - 1320B$ \\
$(0,4)$ & $3265A - 1800B$       & $(2,3)$ & $33743A - 18834B$ \\
$(0,5)$ & $52907A - 29706B$     & $(2,4)$ & $317892B - 569719A$ \\
$(1,1)$ & $12B - 21A$           & $(2,5)$ & $6190602B - 11092229A$ \\
$(1,2)$ & $114B - 205A$         & $(3,3)$ & $295572B - 529549A$ \\
        &                       & $(3,4)$ & $5399634B - 9675553A$ \\
\bottomrule
\end{tabular}
\end{center}
\end{tcolorbox}

\subsubsection*{Duality}

The approximation of the other function of the pair, computed
directly,
\[
  T^-_{mk} = \frac{2^m m!\; 2^k k!}{h_0}
  \int_\R P_m P_k\, \G(\tfrac14+ix)\, \G(\tfrac14-ix)^2\, dx
  \;=\; \overline{T_{mk}}\,,
\]
lies in the same lattice $\Z A + \Z B$. Conjugation flips the
suppressed $i$ to $-i$; at the fixed phase of the $T$-table the
odd-parity entries change sign:

\begin{tcolorbox}[colback=black!5, colframe=black!5, boxrule=0pt, sharp corners]
\begin{center}
\begin{tabular}{cl@{\qquad}cl}
\toprule
$(m,k)$ & $T^-_{mk}$ & $(m,k)$ & $T^-_{mk}$ \\
\midrule
$(0,0)$ & $A$                   & $(1,3)$ & $2575A - 1440B$ \\
$(0,1)$ & $2B - 3A$             & $(1,4)$ & $21978B - 39411A$ \\
$(0,2)$ & $12B - 23A$           & $(1,5)$ & $394932B - 707389A$ \\
$(0,3)$ & $241A - 138B$         & $(2,2)$ & $2369A - 1320B$ \\
$(0,4)$ & $3265A - 1800B$       & $(2,3)$ & $18834B - 33743A$ \\
$(0,5)$ & $29706B - 52907A$     & $(2,4)$ & $317892B - 569719A$ \\
$(1,1)$ & $12B - 21A$           & $(2,5)$ & $11092229A - 6190602B$ \\
$(1,2)$ & $205A - 114B$         & $(3,3)$ & $295572B - 529549A$ \\
        &                       & $(3,4)$ & $9675553A - 5399634B$ \\
\bottomrule
\end{tabular}
\end{center}
\end{tcolorbox}

\subsubsection*{Reverse}

The \emph{reverse} expansion $\widehat P_m = \sum_k D_{mk} F_k$ has
one-Gamma coefficients
\[
  D_{mk} = \int \widehat P_m \widehat P_k\, \G(\tfrac14 - ix)\, dx
\]
over the single generator $e^{-1}$ (a factor of $i$ suppressed when
$m+k$ is odd):
\[
  S_{mk} = 2^m m!\; 2^k k!\; \frac{e}{2\pi}
  \int_\R P_m P_k\, \G(\tfrac14-ix)\, dx \;\in\; \Z :
\]

\begin{tcolorbox}[colback=black!5, colframe=black!5, boxrule=0pt, sharp corners]
\begin{center}
\begin{tabular}{r|rrrrrr}
$S_{mk}$ & $k=0$ & $1$ & $2$ & $3$ & $4$ & $5$ \\
\hline
$m=0$ & $1$ & $3$ & $5$ & $11$ & $249$ & $-2621$ \\
$1$ &  & $7$ & $47$ & $399$ & $-2005$ & $263$ \\
$2$ &  &  & $473$ & $-617$ & $22109$ & $-402065$ \\
$3$ &  &  &  & $31511$ & $-412557$ & $-3017969$ \\
$4$ &  &  &  &  & $-1916623$ & $-29235989$ \\
$5$ &  &  &  &  &  & $-1012892921$ \\
\end{tabular}
\end{center}
\end{tcolorbox}

\begin{tcolorbox}[colback=green!8, colframe=green!8, boxrule=0pt, sharp corners]
\begin{definition}[Defect--chain matching]\label{def:defect-chain}
A \emph{defect--chain matching} is a combinatorial object richer
than a matching. Take two rows of labelled points, $[m]$ and $[k]$.
Within each row, some points are set aside as \emph{defects}; the
remaining points are arranged into a forest of nonempty ordered
\emph{chains}. Finally, chains in the $[m]$-row are partially matched
with chains in the $[k]$-row. Each part carries a weight, and
$S_{mk}$ is the signed sum over all such configurations.
\end{definition}
\end{tcolorbox}

\begin{tcolorbox}[colback=blue!5, colframe=blue!5, boxrule=0pt, sharp corners]
\begin{theorem}[Meixner--Pollaczek $\tfrac14$, reverse direction]\label{thm:mp}
$S_{mk} = (-1)^{\lfloor (m+k)/2 \rfloor} \sum \operatorname{wt}$,
summed over the defect--chain matchings of $[m]$, $[k]$; the
weights are $(-1)^d (2d-1)!!$ per defect set of size $d$,
$(-1)^{\ell-1} 2^{\ell+1}$ per chain of length $\ell$, and $-1$ per
matched pair of chains.

Equivalently, with Lah numbers
$L(n,c) = \binom{n-1}{c-1} n!/c!$,
\[
  S_{mk} = (-1)^{\lfloor \frac{m+k}{2} \rfloor}
  \sum_{d_1, c_1, d_2, c_2}
  W_m(d_1, c_1)\, W_k(d_2, c_2)\,
  \sum_j (-1)^j \tbinom{c_1}{j}\tbinom{c_2}{j} j! \,,
\]
$W_n(d,c) = \binom{n}{d}(-1)^d (2d-1)!!\; L(n{-}d, c)\,
(-1)^{n-d-c}\, 2^{n-d+c}$.
\end{theorem}
\end{tcolorbox}

\subsubsection*{Reverse dual}

The expansion over the other dual function,
$\widehat P_m = \sum_k \widetilde D_{mk}\, G_k$ with
$\widetilde D_{mk} = \int \widehat P_m \widehat P_k\,
\G(\tfrac14 + ix)\, dx$, obeys two symmetries: conjugation
($\widehat P$ is real on $\R$) gives
$\widetilde D_{mk} = \overline{D_{mk}}$, and the substitution
$x \mapsto -x$ gives $\widetilde D_{mk} = (-1)^{m+k} D_{mk}$.
Together they force $D_{mk}$ to be real for $m+k$ even and purely
imaginary for $m+k$ odd --- this is the origin of the suppressed
factor of $i$. Computing directly, at the fixed phase $i$,
\[
  \widetilde S_{mk} = 2^m m!\; 2^k k!\; \frac{e}{2\pi}
  \int_\R P_m P_k\, \G(\tfrac14+ix)\, dx \;\in\; \Z :
\]

\begin{tcolorbox}[colback=black!5, colframe=black!5, boxrule=0pt, sharp corners]
\begin{center}
\begin{tabular}{r|rrrrrr}
$\widetilde S_{mk}$ & $k=0$ & $1$ & $2$ & $3$ & $4$ & $5$ \\
\hline
$m=0$ & $1$ & $-3$ & $5$ & $-11$ & $249$ & $2621$ \\
$1$ &  & $7$ & $-47$ & $399$ & $2005$ & $263$ \\
$2$ &  &  & $473$ & $617$ & $22109$ & $402065$ \\
$3$ &  &  &  & $31511$ & $412557$ & $-3017969$ \\
$4$ &  &  &  &  & $-1916623$ & $29235989$ \\
$5$ &  &  &  &  &  & $-1012892921$ \\
\end{tabular}
\end{center}
\end{tcolorbox}

\begin{tcolorbox}[colback=blue!5, colframe=blue!5, boxrule=0pt, sharp corners]
\begin{theorem}[Meixner--Pollaczek $\tfrac14$, reverse dual]
\label{thm:mpdual}
$\widetilde S_{mk} = (-1)^{m+k}\, S_{mk}
= (-1)^{\lceil (m+k)/2 \rceil} \sum \operatorname{wt}$,
summed over the same defect--chain matchings, with the same weights,
as Theorem~\ref{thm:mp}: only the floor becomes a ceiling.
\end{theorem}
\end{tcolorbox}

%% -------------------------------------------------------------------
\subsection{Meixner--Pollaczek ($\lambda = \tfrac34$): a dual pair}

For the Meixner--Pollaczek polynomials with $\lambda = \tfrac34$ the
kernel is $W = \G(\tfrac34+ix)\,\G(\tfrac34-ix)$ on $\R$, with
$h_n = h_0\, \G(n+\tfrac32)/(n!\,\G(\tfrac32))$,
$h_0 = (\pi/2)^{3/2}$.
We denote the Meixner--Pollaczek polynomials by $Q_n$.
Here $\sqrt W$ has branch points at
$\pm\tfrac{3i}{4}$; the analytic split is $w_\pm = \G(\tfrac34 \pm ix)$
and the dual pair of functions is given by
\[
  F_n = \G(\tfrac34+ix)\,\widehat Q_n, \qquad
  G_n = \G(\tfrac34-ix)\,\widehat Q_n, \qquad
  \int_\R F_m G_n\, dx = \delta_{mn}.
\]
The coefficients of the approximation matrix are given by:
\[
  c^+_{mk} = \int \widehat Q_m \widehat Q_k\,
  \G(\tfrac34+ix)^2\,\G(\tfrac34-ix)\, dx, \qquad
  c^-_{mk} = \int \widehat Q_m \widehat Q_k\,
  \G(\tfrac34+ix)\,\G(\tfrac34-ix)^2\, dx
\]
They are complex conjugates of each other, $c^-_{mk} = \overline{c^+_{mk}}$.
With the rescaling
\[
  \mathcal U_{mk} = \frac{2^{m+k} m!\, k!}{\pi^2/2}
  \operatorname{Re}\bigl[(-i)^{m+k} J_{mk}\bigr],
  \qquad
  J_{mk} = \int Q_m Q_k\, \G(\tfrac34+ix)^2 \G(\tfrac34-ix)\, dx ,
\]
every entry lies in $\Z A' + \Z B'$, where
\[
  A' = U(\tfrac32, 1, 1) \approx 0.35793, \qquad
  B' = U(\tfrac52, 2, 1) \approx 0.21533,
\]
$U$ being Tricomi's confluent hypergeometric function
(no Bessel simplification: $\lambda = \tfrac14$ is the special
Kummer--Bessel locus):

\begin{tcolorbox}[colback=black!5, colframe=black!5, boxrule=0pt, sharp corners]
\begin{center}
\small
\begin{tabular}{cl@{\qquad}cl}
\toprule
$(m,k)$ & $\mathcal U_{mk}$ & $(m,k)$ & $\mathcal U_{mk}$ \\
\midrule
$(0,0)$ & $A'$                    & $(1,3)$ & $7035 A' - 11760 B'$ \\
$(0,1)$ & $3A' - 6B'$             & $(1,4)$ & $128643 A' - 213630 B'$ \\
$(0,2)$ & $39A' - 60B'$           & $(2,2)$ & $6609 A' - 10920 B'$ \\
$(0,3)$ & $513A' - 870B'$         & $(2,3)$ & $112695 A' - 187350 B'$ \\
$(0,4)$ & $8673A' - 14280B'$      & $(2,4)$ & $2215047 A' - 3679860 B'$ \\
$(1,1)$ & $33 A' - 60 B'$         & $(3,3)$ & $2078169 A' - 3456420 B'$ \\
$(1,2)$ & $453 A' - 750 B'$       & $(3,4)$ & $43522881 A' - 72346830 B'$ \\
\bottomrule
\end{tabular}
\end{center}
\end{tcolorbox}

\subsubsection*{Duality}

The approximation of the other function of the pair, computed
directly at the fixed phase of the $\mathcal U$-table,
\[
  \mathcal U^-_{mk} = \frac{2^{m+k} m!\, k!}{\pi^2/2}
  \operatorname{Re}\bigl[(-i)^{m+k}\, \overline{J_{mk}}\,\bigr],
\]
again negates exactly the odd-parity entries; with the phase
conjugated to $i^{m+k}$ the $\mathcal U$-table reappears identically:

\begin{tcolorbox}[colback=black!5, colframe=black!5, boxrule=0pt, sharp corners]
\begin{center}
\small
\begin{tabular}{cl@{\qquad}cl}
\toprule
$(m,k)$ & $\mathcal U^-_{mk}$ & $(m,k)$ & $\mathcal U^-_{mk}$ \\
\midrule
$(0,0)$ & $A'$                    & $(1,3)$ & $7035 A' - 11760 B'$ \\
$(0,1)$ & $6B' - 3A'$             & $(1,4)$ & $213630 B' - 128643 A'$ \\
$(0,2)$ & $39A' - 60B'$           & $(2,2)$ & $6609 A' - 10920 B'$ \\
$(0,3)$ & $870B' - 513A'$         & $(2,3)$ & $187350 B' - 112695 A'$ \\
$(0,4)$ & $8673A' - 14280B'$      & $(2,4)$ & $2215047 A' - 3679860 B'$ \\
$(1,1)$ & $33 A' - 60 B'$         & $(3,3)$ & $2078169 A' - 3456420 B'$ \\
$(1,2)$ & $750 B' - 453 A'$       & $(3,4)$ & $72346830 B' - 43522881 A'$ \\
\bottomrule
\end{tabular}
\end{center}
\end{tcolorbox}

\subsubsection*{Reverse}

The \emph{reverse} expansion $\widehat Q_m = \sum_k D_{mk} F_k$ has
one-Gamma coefficients
\[
  D_{mk} = \int \widehat Q_m \widehat Q_k\, \G(\tfrac34 - ix)\, dx
\]
over the single generator $e^{-1}$ (a phase $(-i)^{|m-k|}$
suppressed):
\[
  S_{mk} = 2^m m!\; 2^k k!\; \frac{e}{2\pi}
  \int Q_m Q_k\, \G(\tfrac34-ix)\, dx \;\in\; \Z :
\]

\begin{tcolorbox}[colback=black!5, colframe=black!5, boxrule=0pt, sharp corners]
\begin{center}
\begin{tabular}{r|rrrrrr}
$S_{mk}$ & $k=0$ & $1$ & $2$ & $3$ & $4$ & $5$ \\
\hline
$m=0$ & $1$ & $-1$ & $-9$ & $-43$ & $-95$ & $1671$ \\
$1$ &  & $15$ & $23$ & $-283$ & $-4767$ & $-52649$ \\
$2$ &  &  & $529$ & $5475$ & $31703$ & $-65119$ \\
$3$ &  &  &  & $-15145$ & $363979$ & $8950653$ \\
$4$ &  &  &  &  & $-8709567$ & $-100346585$ \\
$5$ &  &  &  &  &  & $123082079$ \\
\end{tabular}
\end{center}
\end{tcolorbox}

\begin{tcolorbox}[colback=blue!5, colframe=blue!5, boxrule=0pt, sharp corners]
\begin{theorem}[Meixner--Pollaczek $\tfrac34$, reverse direction]
\label{thm:mp34rev}
$S_{mk} \in \Z$ for all $m,k$, a signed sum over the same
defect--chain matchings as Theorem~\ref{thm:mp}: for $m \le k$,
\[
  S_{mk} = (-1)^{k}
  \sum_{d_1, c_1, d_2, c_2}
  W'_m(d_1, c_1)\, W'_k(d_2, c_2)\,
  \sum_j (-1)^j \tbinom{c_1}{j}\tbinom{c_2}{j} j! \,,
\]
\[
  W'_n(d,c) = \binom{n}{d}(-1)^d (2d+1)!!\; L(n{-}d, c)\,
  (-1)^{n-d-c}\, 2^{n-d+c} .
\]
Only the defect weight changes, from $(2d-1)!!$ to $(2d+1)!!$: the
Charlier kernel $e^{-\sigma-\tau+\sigma\tau}$ is universal in
$\lambda$, and the Jacobian exponent shifts the odd double factorial
by the half-integer $\lambda$.
\end{theorem}
\end{tcolorbox}

\subsubsection*{Reverse dual}

As for $\lambda = \tfrac14$, the expansion over the other dual
function is the complex conjugate; computed directly, at the fixed
phase $(-i)^{|m-k|}$,
\[
  \widetilde S_{mk} = 2^m m!\; 2^k k!\; \frac{e}{2\pi}
  \int Q_m Q_k\, \G(\tfrac34+ix)\, dx \;\in\; \Z :
\]

\begin{tcolorbox}[colback=black!5, colframe=black!5, boxrule=0pt, sharp corners]
\begin{center}
\begin{tabular}{r|rrrrrr}
$\widetilde S_{mk}$ & $k=0$ & $1$ & $2$ & $3$ & $4$ & $5$ \\
\hline
$m=0$ & $1$ & $1$ & $-9$ & $43$ & $-95$ & $-1671$ \\
$1$ &  & $15$ & $-23$ & $-283$ & $4767$ & $-52649$ \\
$2$ &  &  & $529$ & $-5475$ & $31703$ & $65119$ \\
$3$ &  &  &  & $-15145$ & $-363979$ & $8950653$ \\
$4$ &  &  &  &  & $-8709567$ & $100346585$ \\
$5$ &  &  &  &  &  & $123082079$ \\
\end{tabular}
\end{center}
\end{tcolorbox}

\begin{tcolorbox}[colback=blue!5, colframe=blue!5, boxrule=0pt, sharp corners]
\begin{theorem}[Meixner--Pollaczek $\tfrac34$, reverse dual]
\label{thm:mp34dual}
$\widetilde S_{mk} = (-1)^{m+k}\, S_{mk}$: the signed sum of
Theorem~\ref{thm:mp34rev} with the prefactor $(-1)^k$ replaced by
$(-1)^m$ (for $m \le k$), the weights $W'$ unchanged.
\end{theorem}
\end{tcolorbox}

%% ===================================================================
\newpage
\section{Automata}
%% ===================================================================

\begin{tcolorbox}[breakable, colback=yellow!10, colframe=yellow!10, boxrule=0pt, sharp corners]

Each generating function below is a composition: an algebraic
prefactor times $g(\zeta(w,u))$, where $g$ is the kernel of the
family and $\zeta$ an inner function. Extracting the coefficient of
$w^m u^k$ applies $\partial_w$ $m$ times and $\partial_u$ $k$ times,
then evaluates at $w = u = 0$.

Watch the derivatives land on $g(\zeta)$ one at a time:
\[
  \partial_1\, g(\zeta) = g'(\zeta)\,\partial_1\zeta, \qquad
  \partial_2 \partial_1\, g(\zeta)
  = g''(\zeta)\,\partial_2\zeta\,\partial_1\zeta
  + g'(\zeta)\,\partial_2\partial_1\zeta .
\]
At each step the new derivative either hits $g$ again, opening a new
inner factor --- a new \emph{block} --- or deepens an existing one.
Iterating, the terms are indexed by the set partitions $\pi$ of the
$m+k$ derivative slots (Fa\`a di Bruno):
\[
  \partial^{\,m+k} g(\zeta)
  = \sum_{\pi}\; g^{(|\pi|)}(\zeta) \prod_{B \in \pi} \partial_B\, \zeta .
\]
Each term factors in two ways. For each block $B$ the factor
$\partial_B\, \zeta$ is a mixed partial of the \emph{inner} function
at the base point --- a rational constant, essentially a Taylor
coefficient of $\zeta$. So the shape of the partition --- which
slots are grouped with which --- only multiplies the term by a
rational dressing. The outer function is differentiated once per
block, giving $g^{(j)}$ with $j = |\pi|$, and $g^{(j)}$ is what the
automaton computes. Of all the information in $\pi$ the automaton
therefore sees only the single number $j$: two partitions with the
same number of blocks feed it identically, differing only in their
rational dressing.

The walks compute $g^{(j)}$.

When the kernel satisfies a linear ODE
of order $r$, its derivatives close on $r$ states, and each
application of $d/dz$ is an $r \times r$ transfer matrix; expanding
the $j$-fold matrix product over intermediate states is a sum over
walks of length $j$ in the $r$-state automaton, one step per block,
each step weighted by an entry of the transfer matrix.

\medskip
Two refinements: derivatives may also land on the algebraic
prefactor, distributed by the Leibniz rule before the partition is
formed; and the matrix entries are $z$-dependent, so
repeated differentiation also dresses the walk steps with their
derivatives. Both contribute rational factors.

\subsubsection*{A toy example}

Take the Meixner--Pollaczek entry $J_{11}$: two derivative slots on
\[
  G = \frac{\sqrt2\, h_0\; g(\zeta)}{\sqrt{(1-iw)(1-iu)}},
  \qquad g(z) = e^{z/2} K_0(z/2), \quad
  \zeta = \varphi(w)\,\varphi(u).
\]
At the base point $w = u = 0$: $\zeta = 1$ and
\[
  \partial_w \zeta = \partial_u \zeta = 2i, \qquad
  \partial_w \partial_u \zeta = -4,
\]
while a slot landing on the prefactor contributes $i/2$. The walk
states carry the values $\sqrt2\, e^{\zeta/2} K_0(\zeta/2) = A$ and
$\sqrt2\, e^{\zeta/2} K_1(\zeta/2) = B$ at $\zeta = 1$. Each
$d/d\zeta$ steps $K_0 \to K_0$ with weight $\tfrac12$ (the
exponential differentiating) or $K_0 \to K_1$ with weight
$-\tfrac12$; from $K_1$, to $K_0$ with weight $-\tfrac12$ or the
self-loop with weight $\tfrac12 - 1/\zeta = -\tfrac12$ --- the
automaton of the next subsection, dressed by the exponential and
the chain-rule halves. In units of $h_0$:
\begin{center}
\begin{tabular}{llll}
slots & blocks & walk & contributes \\
\hline
both on prefactor & --- & empty walk & $-\tfrac14\, A$ \\
one on $g$ ($\times 2$) & $\{w\}$ or $\{u\}$ & $K_0 \to K_0$ & $-A$ \\
 & & $K_0 \to K_1$ & $+B$ \\
both on $g$, split & $\{w\}\{u\}$ & $K_0 \to K_0 \to K_0$,\; $K_0 \to K_1 \to K_0$ & $-2A$ \\
 & & $K_0 \to K_0 \to K_1$,\; $K_0 \to K_1 \to K_1$ & $0$ \\
both on $g$, merged & $\{wu\}$ & $K_0 \to K_0$ & $-2A$ \\
 & & $K_0 \to K_1$ & $+2B$ \\
\end{tabular}
\end{center}
The split partition walks two steps with Taylor factor
$\partial_w\zeta\,\partial_u\zeta = -4$; the merged block walks one
step and picks up $\partial_w\partial_u\zeta = -4$ instead --- the
Cayley transform makes the two coincide, but they weight walks of
different lengths. Total: $3B - \tfrac{21}{4}A$, i.e.\
$J_{11} = h_0\,(12B - 21A)/4$ --- the entry $T_{11} = 12B - 21A$
of the table in \S3.1.
\end{tcolorbox}

%% -------------------------------------------------------------------
\subsection{Meixner--Pollaczek ($\lambda = \tfrac14$): a dual pair}

The unrescaled coefficients
$J_{mk} = h_0\, T_{mk}/(2^m m!\, 2^k k!)$ have a closed generating
function: with
\[
  \zeta(w,u) = \varphi(w)\,\varphi(u), \qquad
  \varphi(t) = \frac{1+it}{1-it}
  \quad (\text{the Cayley transform}),
\]
\[
  \sum_{m,k \ge 0} J_{mk}\, w^m u^k
  = \frac{2\pi^{3/2}\, e^{\zeta/2}\, K_0(\zeta/2)}
         {\sqrt{(1-iw)(1-iu)}} .
\]

\begin{tcolorbox}[colback=blue!5, colframe=blue!5, boxrule=0pt, sharp corners]
\begin{conjecture}[Meixner--Pollaczek $\tfrac14$, forward direction]\label{conj:mp}
$T_{mk} \in \Z A \oplus \Z B$ for all $m, k$. No explicit
combinatorial formula is known.
\end{conjecture}
\end{tcolorbox}

The automaton has states $\{K_0, K_1\}$ for the kernel
$g(z) = e^{z/2} K_0(z/2)$, with transitions given by the Bessel
relations
\[
  K_0' = -K_1, \qquad K_1' = -K_0 - \frac{K_1}{z},
\]
evaluated at $z = \tfrac12$: a step either switches state (weight
$-1$) or self-loops at $K_1$ (weight $-1/z = -2$). Walks start at
$K_0$; those ending at $K_0$ contribute to the coefficient of $A$,
those ending at $K_1$ to that of $B$.

The same automaton covers the other function of the pair:
conjugating the generating function is the substitution
$(w,u) \mapsto (-w,-u)$ --- the Cayley transform obeys
$\varphi(-t) = \overline{\varphi(t)}$ --- which leaves the kernel,
states and weights unchanged and only flips the sign of the
odd-parity dressings, reproducing the duality of \S3.1.

%% -------------------------------------------------------------------
\subsection{Meixner--Pollaczek ($\lambda = \tfrac34$): a dual pair}

The coefficients $J_{mk}$ have a closed generating function: with
$\zeta$ as above,
\[
  \sum_{m,k \ge 0} J_{mk}\, w^m u^k
  = \frac{\pi^2}{2}\;
    \frac{U\!\left(\tfrac32,\,1,\,\zeta\right)}
         {(1-iw)^{3/2}(1-iu)^{3/2}} .
\]
This is the $\lambda = \tfrac14$ generating function with the exponent
$\tfrac12$ promoted to $\tfrac32$ and the Bessel kernel
$e^{\zeta/2} K_0(\zeta/2)$ replaced by the Tricomi kernel
$U(\tfrac32, 1, \zeta)$.

\begin{tcolorbox}[colback=blue!5, colframe=blue!5, boxrule=0pt, sharp corners]
\begin{conjecture}[Meixner--Pollaczek $\tfrac34$, forward direction]
\label{conj:mp34}
$\mathcal U_{mk} \in \Z A' \oplus \Z B'$ for all $m, k$. No explicit
combinatorial formula is known.
\end{conjecture}
\end{tcolorbox}

The automaton has states $\{U_0, U_1\}$ for the kernel
$U_0 = U(\tfrac32, 1, z)$, $U_1 = U(\tfrac52, 2, z)$, with
transitions given by the Tricomi contiguous relations
\[
  U_0' = -\tfrac32\, U_1, \qquad
  U_1' = -\frac{U_0}{z} + \frac{z-1}{z}\, U_1 ,
\]
evaluated at $z = 1$: the step $U_0 \to U_1$ has weight $-\tfrac32$,
the step $U_1 \to U_0$ has weight $-1$, and the self-loop on $U_1$
has weight $0$. Walks start at $U_0$; those ending at $U_0$
contribute to the coefficient of $A' = U_0(1)$, those ending at
$U_1$ to that of $B' = U_1(1)$. As for $\lambda = \tfrac14$, the
other function of the pair is covered by the same automaton via
$(w,u) \mapsto (-w,-u)$.

%% ===================================================================
\section{The Mellin Transform and the Riemann Zeta Function}
%% ===================================================================

\mbox{}\par\vspace{-1.5\baselineskip}
\begin{tcolorbox}[breakable, colback=yellow!10, colframe=yellow!10, boxrule=0pt, sharp corners]
The Mellin transform sends the Hermite functions to the
Meixner--Pollaczek dual functions.

The Mellin transform is the Fourier transform of the multiplicative
group $\R_+$; on the critical line $s = \tfrac12 + 2ix$ it is
unitary.

\begin{tcolorbox}[colback=black!5, colframe=black!5, boxrule=0pt,
  sharp corners]
\begin{theorem}\label{thm:mellin}
For every $n \ge 0$, with $m = \lfloor n/2 \rfloor$,
\[
  \mathcal M\bigl[\Phi_n(\sqrt2\,\cdot\,)\bigr]
  \bigl(\tfrac12 + 2ix\bigr)
  = \frac{2^{1/4}\sqrt{\pi}}{2}\; i^{m}\, F_m(x),
\]
where $F_m$ belongs to the $\lambda = \tfrac14$ family for $n$ even
and to the $\lambda = \tfrac34$ family for $n$ odd.
\end{theorem}
\end{tcolorbox}

One constant serves every $n$ and both parities. Its modulus is
forced: $|c|^2 = \pi/2\sqrt2$ is the Plancherel constant of the
$\sqrt2$-dilated Mellin transform, so the map is unitary, and $i^m$
is its eigenvalue phase --- the Mellin analogue of the Fourier
eigenvalues $(-i)^n$ of $\Phi_n$. Underneath is a polynomial
identity: the Bump--Ng polynomials (Bump--Ng 1986) at
$s = \tfrac12 + 2ix$ \emph{are} the Meixner--Pollaczek polynomials
(Coffey 2013), and their real-zero property (``local RH'') is
Meixner--Pollaczek orthogonality. The kernel
$\G(\tfrac14+ix)\,\G(\tfrac14-ix)$ is $|\G(s/2)|^2$, the archimedean
factor of the completed zeta function $\xi(s)$; the functional
equation $s \mapsto 1-s$ is $x \mapsto -x$, which swaps
$F \leftrightarrow G$ --- the dual pair \emph{is} the functional
equation, seen locally.

\medskip
One family on $\R$, two on $\R_+$: Mellin lives on the half-line,
where even and odd are no longer related by $x \mapsto -x$, and the
Hermite family splits into two independent halves --- $\G(s/2)$ for
even, $\G(\tfrac{s+1}{2})$ for odd, the archimedean factors of even
and odd Dirichlet $L$-functions. The even Hermite functions alone
are a complete orthogonal system on $\R_+$ (an even function is
determined by its half), and their Mellin image is the complete
system $\{F_n\}$; the biorthogonality
$\int_\R F_m G_n\, dx = \delta_{mn}$ of \S3 is Hermite orthogonality
on $\R_+$ by Plancherel. What is lost is the orthogonality between
the halves: on the half-line $\Phi_{2m}$ and $\Phi_{2n+1}$ are no
longer orthogonal, and their cross-pairing is the unbalanced kernel
$\G(\tfrac14+ix)\,\G(\tfrac34-ix)$, which we treat elsewhere.

\medskip
So the two Meixner--Pollaczek families are not an extension of
Hermite: they \emph{are} the Hermite family, viewed
multiplicatively. The matrices of this paper change basis between a
function and its own Mellin shadow --- and why that produces
integers, and lattices $\Z A \oplus \Z B$, we do not understand.
\end{tcolorbox}

%% ===================================================================
\section{Proofs Available Upon Request}
%% ===================================================================

\end{document}
